% =============================================================================
% RCS Paper 1 --- Empirical Stability Budget (IEEE Conference Format)
% Skeleton scaffold (body to be drafted in P1f).
% =============================================================================
\documentclass[conference]{IEEEtran}
\IEEEoverridecommandlockouts

% --- Encoding and fonts ---
\usepackage[utf8]{inputenc}
\usepackage[T1]{fontenc}

% --- Colors (must load before hyperref) ---
\usepackage{xcolor}

% --- Hyperlinks ---
\usepackage[colorlinks=true,linkcolor=blue!70!black,citecolor=blue!70!black,urlcolor=blue!70!black]{hyperref}

% --- Tables ---
\usepackage{booktabs}
\usepackage{array}

% --- Micro-typography ---
\usepackage{microtype}

\usepackage{cite}

% --- Shared RCS preamble (math, theorems, macros) ---
% \SHAREDLATEX points at the shared LaTeX assets dir. Each paper in
% papers/pN-.../ resolves it relative to its own location; preamble.tex
% uses the macro to reference its own includes (parameters.tex, etc.).
\def\SHAREDLATEX{../../latex}
% =============================================================================
% RCS — Shared LaTeX Preamble
% Used by all papers in the Recursive Controlled Systems series
% =============================================================================

% --- Math ---
\usepackage{amsmath,amssymb,amsthm}
\usepackage{mathtools}
\usepackage{bm}              % Bold math symbols

% --- Theorem environments ---
\theoremstyle{plain}
\newtheorem{theorem}{Theorem}[section]
\newtheorem{lemma}[theorem]{Lemma}
\newtheorem{proposition}[theorem]{Proposition}
\newtheorem{corollary}[theorem]{Corollary}

\theoremstyle{definition}
\newtheorem{definition}[theorem]{Definition}
\newtheorem{example}[theorem]{Example}
\newtheorem{remark}[theorem]{Remark}
\newtheorem{assumption}[theorem]{Assumption}
\newtheorem{law}{Law}

% --- RCS notation macros ---

% Sets and spaces
\newcommand{\X}{\mathcal{X}}           % State space
\newcommand{\Y}{\mathcal{Y}}           % Observation space
\newcommand{\U}{\mathcal{U}}           % Control input space
\newcommand{\M}{\mathcal{M}}           % Mutation operator set
\newcommand{\setR}{\mathbb{R}}         % Real numbers
\newcommand{\setN}{\mathbb{N}}         % Natural numbers
\newcommand{\setZ}{\mathbb{Z}}         % Integers

% RCS 7-tuple
\newcommand{\rcs}{\Sigma}              % A Recursive Controlled System
\newcommand{\plant}{\rcs}              % Alias: plant
\newcommand{\dyn}{f}                   % Dynamics function
\newcommand{\obs}{h}                   % Observation map
\newcommand{\shield}{S}                % Safety shield
\newcommand{\ctrl}{\Pi}                % Controller

% Levels
\newcommand{\Li}[1]{L_{#1}}           % Level i: \Li{0}, \Li{1}, etc.

% Stability budget
\newcommand{\stab}{\lambda}            % Stability margin
\newcommand{\decay}{\gamma}            % Nominal decay rate
\newcommand{\adapt}{L_\theta \rho}     % Adaptation cost
\newcommand{\design}{L_d \eta}         % Design evolution cost
\newcommand{\delay}{\beta \bar{\tau}}  % Delay cost
\newcommand{\switch}{\frac{\ln \nu}{\tau_a}} % Switching cost

% Homeostatic drive
\newcommand{\drive}[1]{D_{#1}}         % Drive at level i: \drive{i}
\newcommand{\setpoint}[1]{x^*_{#1}}   % Setpoint at level i
\newcommand{\homostate}{\xi}           % Homeostatic state vector

% Lyapunov
\newcommand{\lyap}[1]{V_{#1}}         % Lyapunov function at level i
\newcommand{\dlyap}{\dot{V}}          % Time derivative of V

% EGRI
\newcommand{\egri}{\Pi_{\text{EGRI}}}  % EGRI meta-controller
\newcommand{\budget}{B}                % Budget
\newcommand{\score}{J}                 % Evaluator/score function
\newcommand{\trial}{\mathcal{T}}       % Trial record

% Agents and information state
\newcommand{\info}{I}                  % Information state (Eslami)
\newcommand{\mem}{m}                   % Memory state
\newcommand{\toolout}{z}               % Tool output
\newcommand{\interact}{r}              % Interaction signal
\newcommand{\param}{\theta}            % Adaptable parameters
\newcommand{\goal}{\zeta}              % Goal descriptor
\newcommand{\mode}{\sigma}             % Mode/strategy index
\newcommand{\arch}{c}                  % Architecture/workflow config

% Eslami delays
\newcommand{\taubar}{\bar{\tau}}       % Supremal total delay
\newcommand{\tauu}{\tau_u}             % Control computation delay
\newcommand{\tauth}{\tau_\theta}       % Adaptation update delay
\newcommand{\tauz}{\tau_z}             % Tool/retrieval delay
\newcommand{\taus}{\tau_\sigma}        % Switching decision delay
\newcommand{\tauc}{\tau_c}             % Reconfiguration delay
\newcommand{\taua}[1][]{\tau_{a\ifx\\#1\\\else,#1\fi}} % Average dwell time: \taua = tau_a, \taua[i] = tau_{a,i}

% Norms
\newcommand{\norm}[1]{\left\| #1 \right\|}
\newcommand{\abs}[1]{\left| #1 \right|}

% Operators
\DeclareMathOperator*{\argmin}{arg\,min}
\DeclareMathOperator*{\argmax}{arg\,max}


% =============================================================================
\title{Measuring the Recursive Stability Budget\\
In a Production Agent Operating System}

\author{
  \IEEEauthorblockN{Carlos D. Escobar-Valbuena}
  \IEEEauthorblockA{
    Maestr\'ia en Inteligencia Artificial (MAIA)\\
    Universidad de los Andes\\
    Bogot\'a, Colombia\\
    carlosdavidescobar@gmail.com
  }
}

% =============================================================================
\begin{document}
\maketitle

\begin{abstract}
\emph{Scaffold placeholder.} This paper reports a runtime validation of
the Recursive Controlled System stability budget introduced in Paper~0
of this series. Using the Life Agent OS as the experimental plant, we
estimate the per-level margin $\hat{\stab}_i$ and composite rate
$\hat{\omega}_N = \min_i \hat{\stab}_i$ from OpenTelemetry
instrumentation of the autonomic controller. We characterize regime
transitions, budget violations, and the empirical distribution of the
four cost components (adaptation, design, delay, switching) that
Paper~0's theory predicts. The final abstract will be written in
P1f once experimental data has been collected according to
\texttt{PROTOCOL.md}.
\end{abstract}

\begin{IEEEkeywords}
  stability budget, recursive controlled systems, autonomous agents,
  OpenTelemetry, empirical validation, Lyapunov, production monitoring
\end{IEEEkeywords}

% =============================================================================
\section{Introduction}
\label{sec:introduction}
% =============================================================================
% TODO(P1f): motivate empirical validation of the stability budget;
% state the contribution (first production-scale measurement of
% $\hat{\stab}_i$ across L0--L3); position relative to Paper~0's
% inductive theorem.

% =============================================================================
\section{Related Work}
\label{sec:related}
% =============================================================================
% TODO(P1f): Eslami \& Yu 2026 original budget; homeostatic control
% monitoring; observability of switched systems; Lyapunov function
% estimation from data; production ML-safety monitoring.

% =============================================================================
\section{RCS Recap}
\label{sec:recap}
% =============================================================================
% Short, cites Paper~0 of the series: restate the 7-tuple, the
% per-level budget decomposition
% $\stab_i = \decay_i - \adapt_i - \design_i - \delay_i - \switch_i$,
% and the composite rate $\omega_N = \min_i \stab_i$ from
% Paper~0's Theorem thm:recursive. No new proofs here.

% =============================================================================
\section{Experimental Methodology}
\label{sec:methodology}
% =============================================================================
% TODO(P1f): point at PROTOCOL.md for full detail. Sub-sections:
%   A. System under test (Life Agent OS deployment)
%   B. Instrumentation (StabilityBudget, MarginEstimator, vigil gauges)
%   C. Data capture (sampling rate, window length, run matrix)
%   D. Dataset layout (data/p1-runs/)

% =============================================================================
\section{Results}
\label{sec:results}
% =============================================================================
% TODO(P1f): per-level margin traces; composite rate across load
% regimes; budget-violation incidents; cost-component breakdown;
% comparison against the canonical parameters in
% data/parameters.toml.

% =============================================================================
\section{Discussion}
\label{sec:discussion}
% =============================================================================
% TODO(P1f): where the theory fits observations; where it doesn't;
% dominant cost term in practice; implications for controller design
% and for the Paper~0 hypotheses (H1)--(H7).

% =============================================================================
\section{Conclusion}
\label{sec:conclusion}
% =============================================================================
% TODO(P1f): summary; limitations; threads to Papers 2--4.

% =============================================================================
% BIBLIOGRAPHY
% =============================================================================
% Scaffold: anchor at least one entry so the bibliography is non-empty
% until P1f populates the body with real \cite commands. Eslami & Yu's
% original stability budget is the load-bearing reference for Paper 1.
\nocite{eslami2026control}
\bibliographystyle{IEEEtran}
\bibliography{\SHAREDLATEX/references}

\end{document}
